% Architecture diagram for AO-PSO.
% Shows the seven value additions sitting inside the two stages of the
% algorithm (Initialisation and Adaptive Evolution).
% The stage rectangles are rendered on the background layer so they do
% NOT paint over the component nodes. The feedback dashed arrow is
% routed entirely outside the boxes so it does not strike through any
% component.

\begin{tikzpicture}[
    font=\footnotesize,
    node distance=5mm and 6mm,
    stage/.style={
        draw=aopsoBlue, very thick, rounded corners=4pt,
        inner sep=10pt,
    },
    component/.style={
        draw=aopsoBlue, rounded corners=3pt, fill=white,
        minimum width=29mm, minimum height=10mm,
        align=center,
    },
    novel/.style={
        draw=aopsoRed, rounded corners=3pt, fill=aopsoRed!15,
        minimum width=29mm, minimum height=10mm,
        align=center, very thick,
    },
    arrow/.style={-{Latex[length=2mm,width=1.8mm]}, thick},
    fadearrow/.style={-{Latex[length=2mm]}, draw=aopsoGray, dashed, thick},
    label/.style={font=\footnotesize\bfseries, color=aopsoBlue},
]

% ---------- Stage 1: initialisation (top row) ---------------------------
\node[component] (chaos)  {\textbf{VA-1}\\Chaotic Init\\\scriptsize logistic, $\mu{=}4$};
\node[component, right=of chaos] (gobl0)
       {\textbf{VA-2}\\Generalized OBL\\\scriptsize $\hat x = a{+}b{-}\beta x{-}(1{-}\beta)\bar x$};
\node[component, right=of gobl0]  (sel)
       {Select best $N$\\of $2N$ candidates};

% ---------- Stage 2: adaptive evolution (rows below) --------------------
\node[component, below=12mm of chaos] (psoupd)
       {Constriction PSO\\velocity \& move};
\node[component, right=of psoupd] (rebel)
       {\textbf{VA-6}\\Rebel Repulsion\\\scriptsize $-\alpha(p_{\text{worst}}{-}x_i)$};
\node[component, right=of rebel]  (eval)
       {Evaluate swarm\\$f(X)$};

\node[component, below=of psoupd] (sigma)
       {Swarm diameter\\$\sigma_t = \mathrm{std}/\mathrm{width}$};
\node[novel,     right=of sigma]  (adapt)
       {\textbf{VA-4 NOVEL}\\Adaptive $J_r^{(t)}$\\\scriptsize $=J_{r,\max}(1{-}e^{-\lambda\sigma_t})$};
\node[component, right=of adapt]  (jump)
       {\textbf{VA-3} Gen.\ Jumping\\\scriptsize if $U(0,1) < J_r^{(t)}$};

\node[component, below=of sigma] (stag)
       {\textbf{VA-5} pbest opp.\\\scriptsize if $s_i \geq t_{\text{stag}}$};
\node[component, right=of stag]  (upd)
       {Update $p_i$, $g$, $s_i$};
\node[component, right=of upd]   (fe)
       {\textbf{VA-7} FE budget\\\& Wilcoxon test};

% ---------- background stage rectangles --------------------------------
\begin{scope}[on background layer]
  \node[stage, fill=aopsoLight,
        fit=(chaos)(gobl0)(sel)] (init) {};
  \node[stage, fill=aopsoLight!60,
        fit=(psoupd)(rebel)(eval)(sigma)(adapt)(jump)(stag)(upd)(fe)] (evo) {};
\end{scope}

% stage labels
\node[label] at ([yshift=2.5mm]init.north) {Stage 1 -- Initialisation};
\node[label] at ([yshift=2.5mm]evo.north)  {Stage 2 -- Adaptive Evolution Loop};

% ---------- arrows ------------------------------------------------------
\draw[arrow] (chaos) -- (gobl0);
\draw[arrow] (gobl0) -- (sel);

% Stage 1 -> Stage 2 entry: out the bottom of init, into psoupd
\draw[arrow] (sel.south) |- ([yshift=4mm]psoupd.north) -- (psoupd.north);

\draw[arrow] (psoupd) -- (rebel);
\draw[arrow] (rebel)  -- (eval);
\draw[arrow] (eval.south) |- ([yshift=4mm]sigma.north) -- (sigma.north);
\draw[arrow] (sigma) -- (adapt);
\draw[arrow] (adapt) -- (jump);
\draw[arrow] (jump.south) |- ([yshift=4mm]stag.north) -- (stag.north);
\draw[arrow] (stag) -- (upd);
\draw[arrow] (upd)  -- (fe);

% Feedback loop routed COMPLETELY outside the Stage 2 box.
% Goes: fe.south -> down -> all the way left -> up -> back into psoupd.west.
\coordinate (feA) at ([yshift=-7mm]fe.south);
\coordinate (feB) at ([xshift=-7mm,yshift=-7mm]evo.south west);
\coordinate (feC) at ([xshift=-7mm]psoupd.west);
\draw[fadearrow] (fe.south) -- (feA) -- (feB) -- (feC) -- (psoupd.west);
\node[label, color=aopsoGray, font=\scriptsize\itshape]
      at ($(feA)!0.5!(feB)$) [above] {next iteration};

\end{tikzpicture}
